\bookchapter{Lifecycle, Reception and Carrier
Return}\label{viii-ch05-lifecycle-reception-and-carrier-return}

The W8--X1--W9 lifecycle gives the exact program a typed request, response and completion sequence. The corrected carrier construction uses Theta18 for normalization while retaining X1 as the response type. This chapter follows that construction into ordered execution, reception and return, including the controls that distinguish its source roles.

The historical constructions in this chapter retain the fixed 8 September 2026 SAMA source stage. Later results are identified by their own source records.

\subsection{Overview}\label{viii-ch05-overview}

The W8--X1--W9 lifecycle describes one reciprocal transition, not three
freely interchangeable objects. W8 is the unresolved eight-position
substrate state. B/contact prepares the request-side transition. X1 is
the one transient reciprocal response channel. W9 is the resolved
composite state or terminal receipt of the completed relation,

\[
W9=X1+W8.
\]

The equality is typed. It cannot be read as ordinary inventory addition,
as a claim that X1 is matter, or as an equation saying B and X1 are two
scalar halves. The exact program surface can require B and X1 as
differently typed inputs and return W9 while a physical causal weld
remains open. That is not a defect in the lifecycle. It is the reason
the lifecycle is useful: syntax, execution, receipt, carrier and
physical identification can be developed without being collapsed into
one assertion.

The evidence chain is cumulative. CR120V rejected the literal two-half
entity formula while retaining a typed outward/return relay. CR120W then
preserved a failed attempt to weld the A-field carrier directly to X1.
CR120X corrected the route by using Theta18 as a carrier-normalization
quantum while keeping X1 as the response type, and CR120Y stress-tested
that separation. CR120Z and CR120ZB carried the ordered B-then-X1
grammar into exact and finite-shot quantum sandboxes, but neither
installed physical definitions of B or X1.

\subsection{Opening question and conceptual
picture}\label{viii-ch05-opening-question-and-conceptual-picture}

The governing question is:

\begin{quote}
How can an unresolved substrate state receive contact, return reciprocal
information and finish as a resolved write while retaining the identity
of the request, response, carrier and receipt?
\end{quote}

The compact picture is

\begin{itemize}
\tightlist
\item
  W8 unresolved source/site state
\item
  -- B/contact prepares the request --\textgreater{}
\item
  X1 transient reciprocal response
\item
  -- typed resolution --\textgreater{}
\item
  W9 resolved composite state and retained receipt
\end{itemize}

This picture contains two related but nonidentical descriptions.

\begin{enumerate}
\def\labelenumi{\arabic{enumi}.}
\tightlist
\item
  \textbf{State description.} W8 and W9 name the unresolved and resolved
  ends of the relation.
\item
  \textbf{Event description.} B/contact and X1 name the differently
  typed request-side and response-side participation needed to traverse
  between those ends.
\end{enumerate}

The lifecycle therefore keeps endpoint state separate from the event
that produced it. A terminal W9 receipt records that the typed route
completed; it does not turn W9 into a new independent source object.
Similarly, X1 can be required by the type system without the installed
runtime exposing an independently variable X1 payload.

\subsection{Definitions, domains and type
boundaries}\label{viii-ch05-definitions-domains-and-type-boundaries}

\subsubsection{The four lifecycle
objects}\label{viii-ch05-the-four-lifecycle-objects}

\begin{itemize}
\tightlist
\item
  \textbf{W8 / \texttt{S8\_BINARY\_SURFACE}.} Type in this document:
  unresolved binary substrate/site surface; Function: supplies the
  source-side state to be resolved; Explicit nonidentity: not W9; not an
  observed matter row.
\item
  \textbf{B / \texttt{B\_CONTACT\_OPERATOR}.} Type in this document:
  non-scalar contact/request operator; Function: prepares or admits the
  request-side transition; Explicit nonidentity: not X1; not scalar
  \texttt{1/2}.
\item
  \textbf{X1 / \texttt{X1\_AXIS\_SELF\_CHANNEL}.} Type in this document:
  axis/response channel with registered scalar value 1; Function:
  supplies the transient reciprocal response-side occurrence; Explicit
  nonidentity: not Theta18, matter, a row or two half-operators.
\item
  \textbf{W9 / \texttt{W9\_CLOSURE\_WITNESS}.} Type in this document:
  closure witness and resolved composite; Function: records the
  completed typed relation \texttt{X1+W8}; Explicit nonidentity: not a
  third primitive object and not scalar-nine \texttt{CARRIER9}.
\end{itemize}

The base identity is

\[
\operatorname{RESOLVE}(W8,B,X1)\longrightarrow W9,
\qquad W9=\{X1+W8\}.
\]

The braces emphasize a completed relation. The expression is not a
license to remove B, exchange B and X1, or infer that two equal scalar
contributions were added.

\subsubsection{Source row, operator, carrier and receipt are different
domains}\label{viii-ch05-source-row-operator-carrier-and-receipt-are-different-domains}

The lifecycle crosses several domains, each with its own counting rule:

\begin{itemize}
\tightlist
\item
  \textbf{N100 source grammar.} Counted object: 100 independent source
  rows/spins; What is not counted there: five derived W9 receipts, X1,
  non-row A operator.
\item
  \textbf{105-row visible QP lane.} Counted object: all source records
  remain visible; What is not counted there: counted/not-counted overlay
  does not delete a row.
\item
  \textbf{lifecycle runtime.} Counted object: typed operator arguments
  and one result entity; What is not counted there: scalar values are
  not an information-capacity count.
\item
  \textbf{Theta carrier lane.} Counted object: released carrier unit of
  size 18; What is not counted there: X1 response identity and ordinary
  matter lift.
\item
  \textbf{history/readback.} Counted object: terminal and ordered event
  receipts; What is not counted there: an additional material object.
\end{itemize}

The five W9 records in the N100/RH grammar are terminal records of
completed \texttt{X1+W8} relations. They remain visible for custody, but
they contribute no new physical/source row. Counting them as five more
independent objects would change the grammar from 100 to 105 rather than
describe its execution.

\subsubsection{Non-row A and the open causal
seams}\label{viii-ch05-non-row-a-and-the-open-causal-seams}

Volume II supplies a narrower installed provenance chain:

\begin{itemize}
\tightlist
\item
  A\_OPERATOR: established non-row transform
\item
  B/contact -\textgreater{} X1 axis fee -\textgreater{} W9: supported
  typed bridge
\item
  A\_OPERATOR -\textgreater{} B/contact: open
\item
  direct A\_OPERATOR -\textgreater{} X1: open
\end{itemize}

The A-field carrier, \texttt{A\_OPERATOR}, its inactive historical row
proxy, B, X1 and Theta18 are distinct. Participation in one 162-ledger
structure does not fill either open arrow.

\subsubsection{Historical language evidence and current SLC
identity}\label{viii-ch05-historical-language-evidence-and-current-slc-identity}

G06 belongs to the preserved SAM Language v0.6 formal-simulator
evidence. It supplies frozen B prepare/request and X1 response
operators. The later SAM Language v0.7 object is a typed reflective-tape
candidate. Neither language version is the SLC described by this
retained source revision.

The immutable Volume III technical spine records the SLC authority that
was current when that spine was written. At the retained source stage, live authority
superseded that earlier pointer wording: \texttt{SLCQ2-RZ} is the
sole frozen/default SLC, and its byte-identical delegated
arithmetic/application base remains separately typed. This chapter uses
the historical language and Courtroom artifacts as evidence for the
lifecycle; it does not promote their version labels into the current SLC
identity.

\subsection{Typed lifecycle construction from start to
finish}\label{viii-ch05-typed-lifecycle-construction-from-start-to-finish}

\subsubsection{The formal operator
premise}\label{viii-ch05-the-formal-operator-premise}

The language-level premise separates preparation from response. In G06,
B contains a binary flip but is not exhausted by that flip, while X1 has
its own frozen response matrix. The executable checks proceed in four
stages:

\begin{enumerate}
\def\labelenumi{\arabic{enumi}.}
\tightlist
\item
  resolve both operators through the single formal registry;
\item
  verify their frozen hashes and declared input/output shapes;
\item
  apply B to every declared prepare/request column and X1 to every
  response column; and
\item
  apply the involutions and recover the source basis exactly.
\end{enumerate}

The hash, column, mapping and involution checks all pass. Exchanging
request and response, reducing B to its binary-flip suboperation, or
selecting an unfrozen sign/gauge would test a different operator
surface.

\subsubsection{\texorpdfstring{The typed \texttt{RESOLVE}
gate}{The typed RESOLVE gate}}\label{viii-ch05-the-typed-resolve-gate}

CR120A executes

\begin{itemize}
\tightlist
\item
  \nolinkurl{S8_BINARY_SURFACE} : BinarySurface, Layer1StructuralCapacity
\item
  -- RESOLVE(B\_CONTACT\_OPERATOR, \nolinkurl{X1_AXIS_SELF_CHANNEL})
  --\textgreater{}
\item
  \nolinkurl{W9_CLOSURE_WITNESS} : ClosureWitness, Layer2Resolution
\end{itemize}

The runtime does not merely check \texttt{8+1=9}. It enforces the
operator signature. The exact control sequence is:

\begin{itemize}
\tightlist
\item
  RESOLVE(S8, B, X1) -\textgreater{} W9 PASS
\item
  PROMOTE\_VOLUME(W9, D3) -\textgreater{} V27 PASS
\item
  PROMOTE\_VOLUME(S8, D3) TYPE\_REJECT
\item
  RESOLVE(S8, B) ARITY\_REJECT
\item
  RESOLVE(S8, B, B) TYPE\_REJECT
\item
  RESOLVE(S8, X1, X1) TYPE\_REJECT
\item
  W9 == CARRIER9 AUTHORITY\_REJECT
\end{itemize}

This establishes a typed unresolved-to-resolved gate. Its secondary
boundary is equally important: the runtime does not expose an
independently perturbable X1 payload or a recoverable W9 preimage,
filling or correction certificate.

\subsubsection{What the installed runtime computes---and what it only
names}\label{viii-ch05-what-the-installed-runtime-computesand-what-it-only-names}

CR120D opens the runtime and follows the calculation step by step:

\begin{enumerate}
\def\labelenumi{\arabic{enumi}.}
\tightlist
\item
  resolve the three named entities;
\item
  require input types
  \nolinkurl{(BinarySurface,\ ContactOperator,\ AxisChannel)};
\item
  check authority;
\item
  select the registered fixed result entity
  \texttt{W9\_CLOSURE\_WITNESS};
\item
  emit the operator name, inputs, result, typed-plan trace, source trace
  and warning field.
\end{enumerate}

The X1 registry entry is immutable, has scalar value 1 and has no syntax
for preparing \texttt{x0}, preparing \texttt{x1}, resetting X1, or
varying X1 while S8 and B are fixed. The current \texttt{RESOLVE} output
alphabet is the singleton

\[
\{W9\},
\]

so the installed fixed-result runtime contains zero bits of
input-dependent W9 variation. That is a runtime-contract boundary, not a
physical null experiment and not an information-capacity interpretation
of the scalar 9.

\subsubsection{Structural continuation to V27, F81 and
P80}\label{viii-ch05-structural-continuation-to-v27-f81-and-p80}

CR120E executes the research-mode structural route

\begin{itemize}
\tightlist
\item
  S8 --B through X1--\textgreater{} W9
\item
  W9 -\textgreater{} V27 -\textgreater{} F81
\item
  \nolinkurl{RESERVE_CLOSURE_ADDRESS}(F81, X1) -\textgreater{} P80
\end{itemize}

The returned \nolinkurl{P80_PARTICLE_FACE_CONTENT} has value 80 and
authority \texttt{STRUCTURAL\_ONLY}; normal mode rejects the reserve
operation. The supported structural relation is

\begin{itemize}
\tightlist
\item
  P80 structural row inventory
\item
  \begin{itemize}
  \tightlist
  \item
    X1 non-row closure/contact address
  \end{itemize}
\item
  = F81 completed-face capacity.
\end{itemize}

This is narrower than a linear quotient theorem. A rank-one quotient
would require F81 and P80 vector spaces, a linear reserve map and
surjectivity. Those premises are not registered. The source-native
search also found neither a punctured \texttt{9\ensuremath{\times}9} row map nor a
ten-octet grouping. The structural capacity weld therefore remains
separate from a row-to-face decoder.

\subsubsection{Constructed split, residue and mirror-odd
candidate}\label{viii-ch05-constructed-split-residue-and-mirror-odd-candidate}

CR120G asks whether a stateful internal response can be constructed
without losing the residue. Its exact rational candidate satisfies

\begin{itemize}
\tightlist
\item
  parent - left - right - kappa = 0
\item
  omit kappa = e2, scalar residual 1
\item
  kappa\_B(B0 identity) = e2
\item
  kappa\_B(B1 swaps e1/e2) = e1
\item
  H1={[}e1{]} stateful readout = 1
\item
  H2={[}e2{]} stateful readout = 0
\item
  static memoryless control = one fixed value for both histories
\end{itemize}

The mirror model reconstructs both faces, projects its odd sector away
from a predeclared X1 fixture direction and returns rank 80. Transverse
and mixed classes are non-null with exact squared norm 5; balanced,
X1-only and degenerate classes are null. The wrong projector fails to
annihilate X1 and changes the transverse norm from 5 to 4.

This is a validated internal mathematical construction, not an installed
SAM route. No registered F81 vector space, X1 vector embedding, inner
product, loop split or 80-coordinate decoder is supplied by the
language. The model is therefore useful as a discriminating construction
while its physical and installed mapping remains open.

\subsection{Derived W9 receipts and the source
ledger}\label{viii-ch05-derived-w9-receipts-and-the-source-ledger}

\subsubsection{The 105-visible/100-independent
overlay}\label{viii-ch05-the-105-visible100-independent-overlay}

CR120K preserves all 105 visible QP records and marks exactly five
\texttt{p=9,g=0} records as derived W9 witnesses. The
independent-content census is

\begin{itemize}
\tightlist
\item
  visible records 105
\item
  derived W9 witnesses 5
\item
  independent-content records 100
\item
  baseline M\_native 16250.625
\item
  duplicated linear W9 content 50.625 = 405/8
\item
  independent M\_native 16200
\item
  independent mean 162
\end{itemize}

No source record is deleted. The overlay changes only whether the record
is counted as independent content.

\subsubsection{Worked exact parent
examples}\label{viii-ch05-worked-exact-parent-examples}

Each derived p9 record has p8 and p1 parents of the same role and
generation. For example,

\[
M_9(\text{QP093A-0019})=\frac{45}{4},
\]

while

\[
M_8(\text{QP093A-0016})+M_1(\text{QP093A-0001})
=10+\frac54=\frac{45}{4}.
\]

Likewise,

\[
12+\frac32=\frac{27}{2},
\qquad
1+\frac18=\frac98.
\]

All five linear residuals are exactly zero. Isolated parent
perturbations change the derived candidate while the direct-nine
candidate stays fixed, which supplies a dependency discriminator. The
records remain separate because downstream \texttt{qA}, observed-mass
and debit/credit fields can have nonzero residuals even when
\texttt{M\_native} closes linearly.

\subsubsection{The CR120T accounting and binding
boundary}\label{viii-ch05-the-cr120t-accounting-and-binding-boundary}

CR120T carries the ledger proposal through distinct gates rather than
treating one exact total as a universal result. The accounting ladder
closes:

\[
100L=100\cdot162=16{,}200,
\]

\[
100N=100\cdot144=14{,}400,
\qquad
100M=100\cdot126=12{,}600,
\]

with

\[
L=N+\Theta,
\qquad M=N-\Theta,
\qquad N=M+\Theta.
\]

Scalar arithmetic does not select the source route by itself: p6+p3 can
also sum to p9. The typed registry privileges p8+p1 because it supplies
the \texttt{S8}/X1/W9 roles.

The downstream gates do not promote. The best frozen ID-free,
value-blind, shuffle-invariant and conjugate-symmetric F81 rule selects
84 rows with \texttt{M\_native=14077.125}, not the supplied 81 rows with
total 12,600. No three-row exception patch is admitted. The binding
successor replaces one feature with the training-frozen
\texttt{abs(N-Z)/A\^{}(1/3)} term and worsens every reported RMS lane:

\begin{itemize}
\tightlist
\item
  \textbf{CR261 20-row holdout.} Baseline RMS MeV: 3.017113068;
  Candidate RMS MeV: 4.071613571; Change: +1.054500503.
\item
  \textbf{CR277 78-row extension.} Baseline RMS MeV: 7.974144225;
  Candidate RMS MeV: 8.067445504; Change: +0.093301279.
\item
  \textbf{CR277 118-row observed.} Baseline RMS MeV: 6.672483502;
  Candidate RMS MeV: 6.828884165; Change: +0.156400663.
\end{itemize}

The source verdict is
\nolinkurl{BOUNDARY_DISCOVERY_SUPPORTED_PROJECTION_OR_BINDING_OPEN};
the binding lane is \nolinkurl{NO_FROZEN_BINDING_IMPROVEMENT}. The
accounting result remains available without being converted into a
binding coefficient.

\subsection{The preserved deviation
chain}\label{viii-ch05-the-preserved-deviation-chain}

\subsubsection{First deviation: the literal two-half entity
formula}\label{viii-ch05-first-deviation-the-literal-two-half-entity-formula}

CR120V recovers a real reciprocal half-slot grammar. The outward and
inward accounts are

\[
\frac12 SW_{out}+\frac12 WRITE_{out}+\frac12 OUTSIDE=\frac32,
\]

\[
\frac12 INSIDE+\frac12 WRITE_{in}+\frac12 SW_{in}=\frac32,
\]

so the complete paired route is

\[
\frac32+\frac32=3=D.
\]

The failure was assigning this route arithmetic to two entities:

\[
B=\frac12,\qquad X1=\frac12.
\]

B is non-scalar, X1 is an AxisChannel with scalar 1, and the shell has
six half-slots rather than two. CR120V therefore rejects the literal
entity formula while retaining the typed B-side/X1-side relay as an
archival hypothesis. Thirteen of thirteen hashes match and eight of
eight wrong controls reject; no physical intervention is executed.

\subsubsection{Second deviation: candidate uniqueness is not a
weld}\label{viii-ch05-second-deviation-candidate-uniqueness-is-not-a-weld}

CR120W searches 26 partition-one QP occurrences under four frozen gates:

\begin{itemize}
\tightlist
\item
  26 partition-one occurrences
\item
  3 legal nonlocal infrastructure rows
\item
  2 neutral closed carrier-axis shells
\item
  1 non-additive address role
\item
  1 shared environmental-axis survivor
\end{itemize}

The unique survivor is QP093A-0305 \texttt{A\_FIELD\_CARRIER}, with zero
native mass, zero arity, \texttt{a\_kernel\_support} route,
\texttt{carrier\_axis}, non-matter ownership and closed carrier support.
The candidate ledger keeps twelve additive modes at 81, the
mirror/checksum at 81 and closure at 162. Restoring the historical
A-field proxy gives the wrong controls 82 and 163.

The attempted executable statement

\begin{itemize}
\tightlist
\item
  \nolinkurl{REALIZES_AXIS_CHANNEL}(A\_FIELD\_CARRIER, \nolinkurl{X1_AXIS_SELF_CHANNEL})
\end{itemize}

does not close. The result preserves verdict
\nolinkurl{FAIL_CR120W_SOURCE_OR_CANDIDATE_RECONSTRUCTION}, 25/25
frozen source hashes, 11/12 wrong controls and no same-run repair. The
useful discovery is unique typed candidacy; the direct
\texttt{A\_OPERATOR\ -\textgreater{}\ \nolinkurl{X1_AXIS_SELF}\_CHANNEL} edge
remains uninstalled.

\subsubsection{Correction: the Theta-normalized dual-depth
route}\label{viii-ch05-correction-the-theta-normalized-dual-depth-route}

CR120X changes the executable question instead of relabeling the
failure. It keeps A, B, X1 and Theta differently typed and asks whether
an existing carrier quantum normalizes the retained/native/full ledger:

\begin{itemize}
\tightlist
\item
  A route at d=1 R\textsuperscript{(d+1)=R}2=N=144
\item
  B contact 144 -\textgreater{} (126 retained, 18 released)
\item
  Theta normalization 126/18 = 7
\item
  144/18 = 8 = S8
\item
  162/18 = 9 = W9
\item
  normalized release 18/18 = 1 = X1 scalar contact
\item
  typed resolution RESOLVE(S8,B,X1) -\textgreater{} W9
\end{itemize}

Every registered carrier is substituted using native payload and
partition scalar. QP093A-0300 \texttt{TENSOR\_CARRIER}/Theta18 is the
sole candidate that reproduces the complete ladder in both comparisons.
All 32 CR256 rows remain exact to maximum reported residual
\texttt{5.000e-08}; only route depth \texttt{d=1} gives
\texttt{R\^{}(d+1)=144=S8*Theta}.

This correction gives Theta18 the carrier-normalization job and X1 the
response occurrence. Their ratio is one in this construction, but their
types remain different. The campaign records 20/20 source hashes, 11/11
hard gates, 1/1 supportive gate and 20/20 rejected wrong controls. The
canonical ledger is 162; restored-row 163 remains the wrong control.

\subsubsection{Retest: model-wide Theta18
stress}\label{viii-ch05-retest-model-wide-theta18-stress}

CR120Y then removes the possibility that 18 was selected from one
attractive ratio. It asks whether the same carrier survives across
independent evidence:

\begin{itemize}
\tightlist
\item
  \textbf{exact derivation routes.} Result: 6/6 converge on 18.
\item
  \textbf{leave-one-family-out folds.} Result: 6/6 retain Theta18.
\item
  \textbf{structural output family \texttt{\{6,9,18,54\}}.} Result:
  Theta18 alone performs the primary-carrier job.
\item
  \textbf{A/Theta route depth.} Result: \texttt{d=1} only.
\item
  \textbf{QP carrier row.} Result: QP093A-0300 only, in both frozen
  value modes.
\item
  \textbf{Starbreaker ledgers.} Result: 96/96 preserve 162-slot and
  18/126/18 typing.
\item
  \textbf{wrong controls.} Result: 12/12 reject.
\end{itemize}

Theta18 is retained as the current primary substrate-carrier candidate
in the scope of this source result. It remains carrier-only: not X1,
matter, rest mass or a fitted mass term. The retest supplements CR120X;
neither result rewrites CR120V or CR120W.

\subsection{Ordered quantum-sandbox
successors}\label{viii-ch05-ordered-quantum-sandbox-successors}

\subsubsection{Exact minimal ordered
model}\label{viii-ch05-exact-minimal-ordered-model}

CR120Z supplies a prospective minimal realization of the typed order:

\begin{itemize}
\tightlist
\item
  \textbar00\textgreater{} -- B:H\_A -- X1:CNOT\_A-\textgreater B
  --\textgreater{} \textbar Phi+\textgreater{}
\end{itemize}

The resulting Bell state has negativity 0.5, concurrence 1, fidelity 1
and CHSH \texttt{2.82842712475}. B alone, X1 alone, X1-before-B and a
classical correlated-return state all have zero negativity. In the relay
extension, an X1-indexed Bell response on the middle systems heralds a
remote pair with the same negativity and CHSH values, while product
measurement and unresolved false heralds remain separable. Local
marginals stay invariant to maximum trace distance \texttt{1.11e-16}.

The noise thresholds reproduce the analytic boundaries: the first grid
point with nonzero negativity is \texttt{v=0.335}, immediately above
\texttt{1/3}, and the first frozen-setting CHSH violation is
\texttt{v=0.71}, immediately above \texttt{1/sqrt(2)}. Custody closes
10/10 hashes, 9/9 hard gates and 12/12 wrong controls.

This demonstrates quantum capability of the typed order in the sandbox.
H and CNOT are prospective maps, not installed physical definitions of B
and X1; W9 is not used as a Bell witness.

\subsubsection{Finite-shot NV boundary}\label{viii-ch05-finite-shot-nv-boundary}

CR120ZB moves the same order discriminator through three inherited NV
simulator profiles, finite sampling, readout confusion, SPAM,
relaxation, dephasing, drift, linear inversion and PSD reconstruction.
Each reconstructed state uses all 15 nontrivial Pauli observables with
4,096 shots per observable. All 2,160 reconstructed states are physical.

\begin{itemize}
\tightlist
\item
  \textbf{room-temperature NV.} Forward median negativity: 0.4615;
  Forward q05: 0.4496; Maximum control q95: 0.0088; Positive gaps:
  48/48; Median CHSH: 2.6915; CHSH\textgreater2 blocks: 48/48.
\item
  \textbf{cryogenic NV.} Forward median negativity: 0.4708; Forward q05:
  0.4624; Maximum control q95: 0.0123; Positive gaps: 48/48; Median
  CHSH: 2.7337; CHSH\textgreater2 blocks: 48/48.
\item
  \textbf{noisy NV.} Forward median negativity: 0.4426; Forward q05:
  0.4126; Maximum control q95: 0.0104; Positive gaps: 48/48; Median
  CHSH: 2.6256; CHSH\textgreater2 blocks: 48/48.
\end{itemize}

The signal survives and all 12 wrong controls behave correctly, but hard
gates pass 8/9. The exact source verdict is
\nolinkurl{BOUNDARY_X1_ORDERED_QUANTUM_INSTRUMENT_FINITE_SHOT_NV_DISCOVERY__ONE_OR_MORE_FROZEN_GATES_FAILED}.
No live backend ran. The result is a simulator-to-device handoff
boundary, not a physical B/X1 weld or installed SAM operation.

\subsection{What the controls
establish}\label{viii-ch05-what-the-controls-establish}

The controls are part of the derivation, not an appendix to it:

\begin{itemize}
\tightlist
\item
  omitting or exchanging B and X1 verifies that the \texttt{RESOLVE}
  signature is typed rather than scalar arithmetic;
\item
  distinguishing \texttt{W9\_CLOSURE\_WITNESS} from scalar-nine
  \texttt{CARRIER9} prevents a value collision from becoming entity
  identity;
\item
  keeping all 105 records visible while counting 100 prevents source
  deletion from masquerading as an independent-content result;
\item
  perturbing p8 and p1 parents tests dependency rather than only
  equality;
\item
  the 163 restored-row lane shows that reinstalling the historical A
  proxy changes the ledger object;
\item
  the literal-half rejection prevents route slots from being reified as
  B and X1 entities;
\item
  the CR120W 11/12 control outcome prevents unique candidacy from being
  promoted through a failed reconstruction;
\item
  leave-one-family-out replay tests whether Theta18 survives removal of
  any one evidence family;
\item
  reversed-order, product, classical and false-herald quantum controls
  show that ordinary correlation or label balance is not the ordered
  signal; and
\item
  the 8/9 hard-gate result in CR120ZB keeps a strong simulated signal
  and an unpassed device-handoff gate visible at the same time.
\end{itemize}

\subsection{Established result and current
boundaries}\label{viii-ch05-established-result-and-current-boundaries}

The source-bound lifecycle can now be stated compactly:

\begin{enumerate}
\def\labelenumi{\arabic{enumi}.}
\tightlist
\item
  W8, B, X1 and W9 have distinct registered types.
\item
  The formal runtime requires \texttt{RESOLVE(S8,B,X1)} and returns W9.
\item
  W9 is a resolved composite receipt, not a new source row or object.
\item
  The fixed runtime does not expose independent X1 intervention or an
  input-dependent W9 outcome alphabet.
\item
  The literal B-half plus X1-half entity equation is rejected; the typed
  outward/return relay remains useful.
\item
  QP093A-0305 is the unique typed A-field-carrier candidate in CR120W,
  while the direct executable A-to-X1 weld remains open.
\item
  Theta18, not X1, performs the tested carrier-normalization job across
  the 7\textbar8\textbar9 ladder and the model-wide stress surface.
\item
  Exact and finite-shot quantum sandboxes preserve the ordered B-then-X1
  discriminator, but no hardware or physical weld is installed.
\item
  CR120T closes ledger accounting while leaving the F81 projection and
  binding transfer open.
\item
  None of these historical test labels changes the live
  \texttt{SLCQ2-RZ} current identity.
\end{enumerate}

No atomic record assigned to this document carries an authorized project
result classification. This revision therefore preserves exact source
statuses and verdicts without manufacturing one.

\subsection{Connections and forward
handoff}\label{viii-ch05-connections-and-forward-handoff}

\begin{itemize}
\tightlist
\item
  \textbf{\texttt{SAMA-D000025}.} Connection: Supplies the Volume II
  connector, substrate-atom and container types and the open A-to-B
  seam..
\item
  \textbf{\texttt{SAMA-D000034}.} Connection: Uses the five derived W9
  receipts as a non-counted boundary in the N100 census..
\item
  \textbf{\texttt{SAMA-D000035}.} Connection: Owns the parser, type
  system, operator registry and state/history semantics underlying the
  formal lifecycle evidence..
\item
  \textbf{\texttt{SAMA-D000040}.} Connection: Develops Theta18 as
  carrier, coupled address and M126-to-N144-to-L162 return..
\item
  \textbf{\texttt{SAMA-D000042}.} Connection: Carries the resolved
  lifecycle into directional history and event-reel computation..
\item
  \textbf{\texttt{SAMA-D000044}.} Connection: Uses the typed lifecycle
  in the cross-scale Nonzero Amplification bridge..
\item
  \textbf{\texttt{SAMA-D000045}.} Connection: Uses retained W8/X1/W9
  history and Theta activity in relational composition..
\end{itemize}

The immediate forward handoff is D40: keep X1 as the response channel,
move the carrier burden to Theta18, derive its coupled address, and
follow the retained state through M126, N144 and L162 without turning a
readout ceiling into independent matter.

\subsection{Sources and evidential
scope}\label{viii-ch05-sources-and-evidential-scope}

This chapter is assembled from the following fixed source manuscripts.
Source-reported campaign outcomes retain their original scope; no new
campaign was executed for this edition.

\begin{itemize}
\tightlist
\item
  \href{https://github.com/iwtbotiwtwot/SAMA/blob/4c2b5dc1649da5900c90b90d0ce56c223dbb7717/vol_iii/03-writes-and-history/w8-x1-w9-lifecycle.md}{W8
  x1 w9 lifecycle}.
\end{itemize}
